\begin{frame}{LCP를 매 스텝 푸는 것의 값}
  \begin{enumerate}
    \item 공이 바닥에 부딪혀 튕기거나 멈추는 행위가 \textbf{계산이
      깨지지 않고}(폭주/NaN 없이) 자연스러워짐
    \item 다리 로봇의 발이 바닥에 붙었다 떨어지기를 반복하는
      \textbf{걸음} 같은 복잡한 접촉도 같은 식으로 처리
  \end{enumerate}
  \vskip0.6em
  \begin{block}{반대편: ``단단한 막대'' 방식}
    접촉을 강제 제약으로만 푸는 엔진: \textbf{여러 접촉이 동시에}
    일어나는 경우(예: 네 다리로 서 있는 로봇) 제약이 서로
    충돌해 \textbf{계산이 불안정}해지기 쉬움
  \end{block}
  \vskip0.5em
  \begin{itemize}
    \item 14.2\,·\,14.3에서 ``다리 로봇을 학습시킬 땐 왜 정밀 물리엔진이
      필요한가''를 논할 때, 그 필요성의 \textbf{뿌리}가 바로 이
      LCP 접촉 솔버에 있다
  \end{itemize}
\end{frame}
